\newpage
	
\section{Первая лекция. 07.09.26. Введение в теорию множеств. Общие понятия.}
\text{Литература для курса:}
\begin{enumerate}
	\item Верещагин, Шень - "Начала теории множеств"
	\item Зорич - "Математический анализ" (I том) (\textit{Обратить внимание, рекоммендуется как объёмный учебный курс})
	\item Фихтенгольц - "Курс дифференциального и интегрального исчисления"
	\item Рудин - "Основы математического анализа"
	\item Теренс Тао - ""
	\item Макаров, Подкорытов, et.al. - "Избранные задачи по вещественному анализу" (\textit{Обратить внимание, рекоммендуется как задачник})
	\item Демидович - "Сборник задач и упражнений по математическому анализу" (\textit{Классический задачник}) См. также т.н. "Антидемидович"
\end{enumerate}
	
\subsection{Принцип индукции.}
Пусть $P(n)_{n\in \mathbb{N}}$ - набор утверждений, индексированный натуральными числами. Пусть $P_0$ - истинно (\textit{база
индукции}), а также пусть $P_k \text{ - истинно} \implies P_{k+1} \text{ - истинно (\textit{предположение индукции}).}$ Тогда $P(n) \text{ -} \\ \text{истинно } \forall n \in\mathbb{N}$.

\begin{proposition}
	Следующие условия эквивалентны:
	\text{1) принцип индукции и 2)} $ \forall M \subseteq \N; M \neq \varnothing; \text{ содержит} \\
	 \text{ наименьший элемент.} $
\end{proposition}

\begin{proof}
	$\text{Докажем сначала, что 2)} \implies \text{1). Пусть } P(0) \text{ - истинно; выполнено **, но }  P(n) \text{ выполнено не} \\ \text{для всех } n \in \N. \text{ Рассмотрим } A \subseteq \N - \text{ те натуральные числа, для которых } P(n) \text{ - ложно, } A \neq \varnothing.$
\end{proof}

	
\subsection{Отображения}


